% ================================================================
% ==== FILE: content/m_spaces/m12.tex
% ================================================================

% ==============================
%  Ontology of Continua — Core
%  Meta-Space: M12
%  FULL MODULE — FINAL VERSION
% ==============================

\subsubsection{$M_{12}$ Обзор}
\label{sec:m12-overview}

$M_{12}$ — высшее мета-пространство, необходимое для континуумов до уровня $K_{12}$.
В то время как $M_{11}$ содержит структурированные башни мета-теоретических иерархий,
$M_{12}$ содержит \emph{универсальные мета-логические среды}, управляющие
допустимостью, когерентностью и глобальными ограничениями для целых \emph{семейств
мета-иерархий}.

Кратко:
\[
M_{12} = \text{минимальное мета-логическое пространство, в котором вся башня }
M_1, \ldots, M_{11} \text{ может существовать согласованно.}
\]

Это не просто «уровень выше»;
здесь вводится новый класс различий (Ax. 21),
новая ось, недоступная для $M_{11}$,
и глобальное ограничение, управляющее всеми предыдущими M-пространствами.

Это делает $M_{12}$ терминальным окружением для Core~1.3.

% ---------------------------------------------------------------
\subsubsection*{1. Допустимое пространство конфигураций \texorpdfstring{$\Omega(M_{12})$}{\Omega(M_{12})}}

\[
\Omega(M_{12}) =
\left\{
\left( \mathfrak{U}, A_{12}, \Theta_{12},
C_{12}, \Lambda_{12}, T_{12} \right)
\ \middle|\
\text{выполняются условия глобальной мета-логической допустимости}
\right\}.
\]

Где:

\begin{itemize}
    \item $\mathfrak{U}$ — допустимые \emph{вселенные мета-иерархий},
          т.е. коллекции башен из $M_{11}$ с совместимыми отношениями;
    \item $A_{12}$ — оси, определяющие вариативность вселенных мета-иерархий,
    \item $\Theta_{12}$ — глобальные пороги, обеспечивающие логическую допустимость,
    \item $C_{12}$ — циклы, определяющие стабилизацию на универсальном уровне,
    \item $\Lambda_{12}$ — универсальные ограничения (логические, структурные, категориальные),
    \item $T_{12}$ — глобальное структурное натяжение между семьями иерархий.
\end{itemize}

Допустимость требует:

\begin{enumerate}
    \item когерентность всех вложенных мета-иерархий (\(M_1\)–\(M_{11}\)),
    \item существование хотя бы одного универсального стабилизирующего цикла,
    \item глобальная согласованность с оператором $\Psi$ (генерация мета-пространства),
    \item ограниченность мета-тензора натяжения $T_{12}<\infty$,
    \item сохранение универсальных структурных законов $\Lambda_{12}$.
\end{enumerate}

По Теореме 8:

\[
\Omega(K_{12}) \subseteq \Omega(M_{12})
\quad \text{необходимо и достаточно для существования }K_{12}.
\]

% ---------------------------------------------------------------
\subsubsection*{2. Оси $A(M_{12})$}

\[
A(M_{12}) =
\{
A_{\mathrm{universal}},\
A_{\mathrm{meta\mbox{-}logical}},\
A_{\mathrm{cross\mbox{-}hierarchy}},\
A_{\mathrm{invariance}},\
A_{\mathrm{global\ compatibility}},\
A_{\mathrm{admissibility}}
\}.
\]

Где:

\begin{itemize}
    \item $A_{\mathrm{universal}}$ — вариативность по всем вселенным мета-иерархий,
    \item $A_{\mathrm{meta\mbox{-}logical}}$ — глобальные логические ограничения,
    \item $A_{\mathrm{cross\mbox{-}hierarchy}}$ — вариативность по структурам $M_i$,
    \item $A_{\mathrm{invariance}}$ — инварианты, сохраняемые на всех уровнях,
    \item $A_{\mathrm{global\ compatibility}}$ — ось совместимости для полных стеков $K/M$,
    \item $A_{\mathrm{admissibility}}$ — ось существования самого $K_{12}$.
\end{itemize}

Это вводит новый класс различий, который нельзя вложить в $M_{11}$:

\[
A_{\mathrm{universal}}\notin A(M_{11}),
\]
следовательно по Ax. 21:

\[
\dim(M_{12}) > \dim(M_{11}).
\]

% ---------------------------------------------------------------
\subsubsection*{3. Потенциалы $P(M_{12})$}

\[
P(M_{12}) =
\{
U_{\mathrm{universal}},\
U_{\mathrm{coherence}},\
U_{\mathrm{global\ inference}},\
U_{\mathrm{meta\mbox{-}stability}},\
U_{\mathrm{admissibility}},\
U_{\mathrm{invariant\ preservation}}
\}.
\]

Интерпретация:

\begin{itemize}
    \item $U_{\mathrm{universal}}$ — потенциал, поддерживающий целые вселенные мета-иерархий,
    \item $U_{\mathrm{coherence}}$ — способность поддерживать когерентность между всеми встроенными M-пространствами,
    \item $U_{\mathrm{global\ inference}}$ — потенциал для рассуждений между вселенными,
    \item $U_{\mathrm{meta\mbox{-}stability}}$ — способность сопротивляться коллапсу целых башен,
    \item $U_{\mathrm{admissibility}}$ — обеспечивает $\Omega(K_{12})\subseteq\Omega(M_{12})$,
    \item $U_{\mathrm{invariant\ preservation}}$ — обеспечивает инварианты между уровнями.
\end{itemize}

$K_{12}$ существует тогда и только тогда, когда:

\[
U_{\mathrm{admissibility}} > \Theta_{12}^{\mathrm{existence}}.
\]

% ---------------------------------------------------------------
\subsubsection*{4. Пороги $\Theta(M_{12})$}

Критические пороги:

\[
\Theta_{12}=
\{
\Theta_{12}^{\mathrm{coherence}},\
\Theta_{12}^{\mathrm{universal}},\
\Theta_{12}^{\mathrm{collapse}},\
\Theta_{12}^{\mathrm{invariance}},\
\Theta_{12}^{\mathrm{reflection}}
\}.
\]

Интерпретация:

\begin{itemize}
    \item $\Theta_{12}^{\mathrm{coherence}}$ — минимальная когерентность среди всех $M_i$,
    \item $\Theta_{12}^{\mathrm{universal}}$ — минимальная универсальная совместимость,
    \item $\Theta_{12}^{\mathrm{collapse}}$ — порог, на котором целые мета-иерархии дезинтегрируются,
    \item $\Theta_{12}^{\mathrm{invariance}}$ — минимальное сохранение инвариантов,
    \item $\Theta_{12}^{\mathrm{reflection}}$ — порог самосогласованности мета-логики.
\end{itemize}

Превышение $\Theta_{12}^{\mathrm{collapse}}$ даёт:

\[
\mathfrak{U} \to \{ M_i \}_{i=1}^{11}
\quad\text{с потерей универсальной структуры}.
\]

% ---------------------------------------------------------------
\subsubsection*{5. Потоки $J(M_{12})$}

\[
J(M_{12})=
(
J_{\mathrm{universal}},\
J_{\mathrm{coherence}},\
J_{\mathrm{reflection}},\
J_{\mathrm{meta\mbox{-}inference}},\
J_{\mathrm{cross\mbox{-}hierarchy}},\
\nabla T_{12},\
J_{\mathrm{stability}}
).
\]

Интерпретация:

\begin{itemize}
    \item $J_{\mathrm{universal}}$ — потоки, переорганизующие вселенные мета-иерархий,
    \item $J_{\mathrm{coherence}}$ — потоки, восстанавливающие когерентность между $M_1$–$M_{11}$,
    \item $J_{\mathrm{reflection}}$ — самореферентные мета-логические потоки,
    \item $J_{\mathrm{meta\mbox{-}inference}}$ — вывод на всём множестве вселенных,
    \item $J_{\mathrm{cross\mbox{-}hierarchy}}$ — стабилизация на разных уровнях $M$,
    \item $\nabla T_{12}$ — градиенты глобального структурного натяжения,
    \item $J_{\mathrm{stability}}$ — потоки, предотвращающие коллапс универсальной структуры.
\end{itemize}

% ---------------------------------------------------------------
\subsubsection*{6. Циклы $C(M_{12})$}

Фундаментальные циклы:

\begin{itemize}
    \item $C_{\mathrm{universal}}$ — образование и стабилизация вселенных мета-иерархий,
    \item $C_{\mathrm{coherence}}$ — петли восстановления когерентности между всеми M-пространствами,
    \item $C_{\mathrm{reflection}}$ — рекурсивные циклы, поддерживающие мета-логическую целостность,
    \item $C_{\mathrm{invariance}}$ — циклы, обеспечивающие глобальные инварианты,
    \item $C_{\mathrm{admissibility}}$ — циклы, обеспечивающие существование $K_{12}$.
\end{itemize}

Критерий жизни:

\[
\oint_{C_{\mathrm{universal}}} dA_{\mathrm{universal}} \approx 0.
\]

% ---------------------------------------------------------------
\subsubsection*{7. Граница $\partial\Omega(M_{12})$}

При приближении к границе:

\begin{itemize}
    \item универсальная несовместимость,
    \item нарушение глобальных структурных инвариантов,
    \item мета-логическая некогерентность,
    \item коллапс стабилизирующих циклов,
    \item расходимость мета-тензора натяжения.
\end{itemize}

Пересечение границы приводит к редукции к изолированным вселенным $M_{11}$.

% ---------------------------------------------------------------
\subsubsection*{8. Связь с $M_{11}$ и полной иерархией K/M}

\[
M_{11} \subset M_{12}
\quad\text{строго}.
\]

Проекция:

\[
\pi_{11}: M_{12} \to M_{11}
\]
сбрасывает универсальные ограничения и оставляет только иерархически структурированные башни.

$M_{12}$ — минимальное окружение, в котором:

\[
\Psi: M_{11} \to M_{12}
\]
определено (Run 11).

Существование $K_{12}$ требует:

\[
\Omega(K_{12}) \subseteq \Omega(M_{12}).
\]

$M_{12}$ — терминальное мета-пространство для Core.
